\documentclass[12pt]{article}
\usepackage[margin=1in]{geometry}
\usepackage{amsmath, amssymb, mathtools}
\usepackage{lmodern}
\usepackage[T1]{fontenc}
\usepackage{microtype}
\usepackage{enumitem}

\setlist[enumerate]{itemsep=0.5em}
\setlength{\parskip}{0.6em}
\setlength{\parindent}{0pt}

\begin{document}

% ============================ QUESTIONS ============================
\section*{Posterior for Poisson Likelihood with Gamma Prior}

\textbf{Goal.} Derive the posterior distribution for the Poisson rate parameter $\lambda$ for a Poisson likelihood and Gamma prior on $\lambda$.

\textbf{Likelihood.} We observe $N$ independent counts $x = \{x_1, x_2, \dots, x_N\}$, and we model
\[
X_i \mid \lambda \;\sim\; \text{Poisson}(\lambda)\quad\text{independently for }i=1,\dots,N.
\]
The Poisson pmf is
\[
\Pr(X_i = x_i \mid \lambda) \;=\; \frac{\lambda^{x_i} e^{-\lambda}}{x_i!},\qquad x_i \in \{0,1,2,\dots\},\ \lambda>0.
\]
Under independence, the \emph{likelihood} of $\lambda$ given the data is
\[
f(x \mid \lambda)\;=\;\prod_{i=1}^N \frac{\lambda^{x_i} e^{-\lambda}}{x_i!}
\;=\; \left[\frac{1}{\prod_{i=1}^N x_i!}\right] e^{-N\lambda}\,\lambda^{\sum_{i=1}^N x_i}.
\]

\textbf{Prior.} Place a Gamma prior on $\lambda$ with shape $\alpha^*>0$ and rate $\beta^*>0$.
\[
\lambda \;\sim\; \text{Gamma}(\alpha^*, \beta^*),\qquad
f(\lambda) \;=\; \frac{(\beta^*)^{\alpha^*}}{\Gamma(\alpha^*)}\,\lambda^{\alpha^*-1} e^{-\beta^* \lambda},\quad \lambda>0.
\]

\textbf{Task.} Find the posterior distribution $f(\lambda \mid x)$, identifying it by name and parameters. (Hint: \textit{The Gamma distribution is conjugate to the Poisson likelihood.})

\newpage

\section*{Explaining the Rejection Algorithm}

\textbf{Task.} In your own words, explain the logic of the \emph{rejection algorithm} to someone who has never seen it before. Assume:
\begin{itemize}
\item The target distribution has support on $[0,1]$ and is known up to a proportionality constant; denote its (possibly unnormalized) density by $f(z)$ for $z \in [0,1]$.
\item We use a \emph{uniform} proposal on $[0,1]$, i.e., propose $Z \sim \text{Uniform}(0,1)$.
\item We have an \emph{envelope constant} $M>0$ satisfying $M \ge \sup_{z\in[0,1]} f(z)$.
\end{itemize}
Explain why the algorithm works and what the output represents. You may use sentences or bullet points, but explain with words rather than equations. And explain \textit{why} it works rather than what it does.

\newpage

% ============================ SOLUTIONS ============================
\begin{center}
{\Large Solutions}\\[2pt]
\end{center}

\section*{Solution to Question 1}

\textbf{Step 1: Write down the pieces explicitly.}

From the pmf and independence,
\[
\begin{aligned}
f(x \mid \lambda)
&=\prod_{i=1}^N \frac{\lambda^{x_i} e^{-\lambda}}{x_i!}
= \left[\frac{1}{\prod_{i=1}^N x_i!}\right]
\left(\prod_{i=1}^N \lambda^{x_i}\right)
\left(\prod_{i=1}^N e^{-\lambda}\right) \\
&= \underbrace{\left[\frac{1}{\prod_{i=1}^N x_i!}\right]}_{\text{constant in }\lambda}\;
\lambda^{\sum_{i=1}^N x_i}\; e^{-N\lambda}.
\end{aligned}
\]
The Gamma prior density (shape--rate) is
\[
f(\lambda)=\underbrace{\frac{(\beta^*)^{\alpha^*}}{\Gamma(\alpha^*)}}_{\text{constant in }\lambda}\;
\lambda^{\alpha^*-1} e^{-\beta^* \lambda},\qquad \lambda>0.
\]

\textbf{Step 2: Multiply likelihood and prior (Bayes up to proportionality).}

Ignoring constants that do not depend on $\lambda$,
\[
\begin{aligned}
f(\lambda \mid x) &\propto f(x\mid \lambda)\,f(\lambda) \\
&\propto \Big(\lambda^{\sum_{i=1}^N x_i} e^{-N\lambda}\Big)\;
\Big(\lambda^{\alpha^*-1} e^{-\beta^* \lambda}\Big).
\end{aligned}
\]

\textbf{Step 3: Collect like terms carefully (same base $\lambda$ and same base $e$).}

For powers of $\lambda$:
\[
\lambda^{\sum_{i=1}^N x_i}\times \lambda^{\alpha^*-1}
= \lambda^{\left(\sum_{i=1}^N x_i + \alpha^* - 1\right)}.
\]
For exponentials:
\[
e^{-N\lambda}\times e^{-\beta^*\lambda}=e^{-(N+\beta^*)\lambda}.
\]
Therefore,
\[
f(\lambda \mid x) \propto \lambda^{\left(\alpha^*+\sum_{i=1}^N x_i\right)-1}\; e^{-(\beta^*+N)\lambda}.
\]

\textbf{Step 4: Name the updated parameters and recognize the family.}

Define
\[
\alpha' = \alpha^* + \sum_{i=1}^N x_i,\qquad \beta' = \beta^* + N.
\]
Then
\[
f(\lambda \mid x) \propto \lambda^{\alpha'-1} e^{-\beta' \lambda},
\]
which is the \emph{kernel} of a $\text{Gamma}(\alpha',\beta')$ density.

\textbf{Step 5: Restore the normalizing constant (to make it a proper density).}

The normalized Gamma density with shape $\alpha'$ and rate $\beta'$ is
\[
f(\lambda \mid x)=\frac{(\beta')^{\alpha'}}{\Gamma(\alpha')}\;\lambda^{\alpha'-1} e^{-\beta'\lambda},\qquad \lambda>0.
\]
Hence the posterior is
\[
\boxed{\ \lambda \mid x \sim \text{Gamma}\!\left(\alpha^* + \sum_{i=1}^N x_i,\; \beta^* + N\right). \ }
\]

\textbf{Why this makes sense (quick checks).}
\begin{itemize}
\item \emph{Conjugacy:} Gamma prior + Poisson likelihood $\Rightarrow$ Gamma posterior.
\item \emph{Learning from data:} More events (large $\sum x_i$) increase the posterior shape; more observations ($N$) increase the rate, shrinking toward stable estimates.
\end{itemize}

\section*{Solution to Question 2}


\textbf{Paragraph explanation}  
Think of the target distribution as a wavy “height” curve drawn above the interval $[0,1]$, and imagine a flat roof set high enough to sit above that entire curve. We first spread candidate points evenly across the floor from 0 to 1 (uniform proposals), which means every location is tried equally often. Then we keep each candidate with a chance that depends only on how tall the target curve is at that spot relative to the roof: where the curve is tall, candidates are kept often; where it is short, they are kept rarely. 

This “shape-based thinning” turns an even carpet of proposals into a pile of kept points whose density mirrors the curve itself. The unknown normalizing constant of the target never matters, because we only compare relative heights; stretching or shrinking the entire curve by a common factor just changes how many proposals are kept overall, not \emph{where} they are kept. The roof guarantees the keep–discard chance is valid everywhere and also determines efficiency (a lower roof, just above the peak, wastes fewer proposals). 

The accepted $z$’s are thus independent draws that behave as if they were generated directly from the target distribution, ready for summaries like means, quantiles, and credible intervals.

\end{document}
